\chapter{Bereavement}

\section*{Definition and Overview}
Bereavement is the state of having lost a loved one through death, and the grief that follows is the normal and natural psychological response to that loss. Grief is not an illness: it is a universal human experience that encompasses emotional, physical, cognitive, and behavioural reactions, which typically come in waves and gradually become less intense over time.

Although most people adjust to bereavement without professional treatment, the grieving process is different for everyone and there is no fixed timetable. In a minority of people grief becomes prolonged and severely impairing, a condition now formally recognised as prolonged grief disorder, for which specific psychological treatment is effective.

\section*{Epidemiology}
Virtually every person experiences bereavement at some point in life, making it a universal human experience.

\begin{itemize}
    \item \textbf{Adjustment:} most people adapt to their loss within one to two years
    \item \textbf{Prolonged grief disorder:} roughly 1 in 10 bereaved adults develops a form of grief in which severe symptoms persist beyond 12 months and interfere with daily functioning
    \item \textbf{High-risk circumstances:} the death of a child or partner, sudden or violent deaths, and multiple losses close together increase the risk of complicated grief
    \item \textbf{Age:} older people experience multiple losses over a lifetime and may face accumulating grief
    \item \textbf{Sex:} men and women are affected equally, although they may express and process grief differently
\end{itemize}

\section*{Aetiology and Pathophysiology}
Grief is a biological and psychological response to the disruption of a strong attachment bond. When a close relationship ends in death, the brain circuits that underpin attachment, reward, memory, and stress regulation are suddenly deprived of their usual input.

In the early months, brain regions involved in emotional bonding and recollection remain highly active, which is why the deceased is so often in the person's thoughts and why ordinary reminders can trigger intense reactions. Stress hormone systems are activated, which can disturb sleep, appetite, and immune function. Over time the brain normally integrates the loss, so that memories no longer trigger overwhelming distress. In prolonged grief this adaptation appears to become stuck: the person remains intensely preoccupied with the loss and cannot re-engage with ongoing life. Vulnerability is influenced by the circumstances of the death, the nature of the relationship, past losses, social support, and a personal or family history of anxiety or depression.

\section*{Clinical Features}
Grief affects the whole person, and the features vary widely between individuals and change over time.

\begin{itemize}
    \item \textbf{Emotional symptoms:} sadness, yearning and longing, emptiness, anger, guilt, anxiety, numbness, and sometimes relief, particularly after a long or difficult illness
    \item \textbf{Physical symptoms:} fatigue, disrupted sleep, poor appetite, a heaviness in the chest, muscle aches, and increased susceptibility to infection
    \item \textbf{Cognitive symptoms:} preoccupation with the deceased, difficulty concentrating, forgetfulness, and vivid memories or dreams
    \item \textbf{Behavioural symptoms:} crying, social withdrawal, and either avoiding or actively seeking reminders of the person
    \item \textbf{Intensity and duration:} symptoms commonly come in waves, often triggered by anniversaries or reminders, and gradually soften over months
    \item \textbf{Prolonged grief disorder:} intense longing and preoccupation persisting beyond 12 months, with severe impairment in work, family, and social life
\end{itemize}

\section*{Differential Diagnosis}
Grief can resemble other conditions, and distinguishing between them matters because the treatments differ.

\begin{itemize}
    \item \textbf{Major depressive disorder:} characterised by persistent low mood across most situations, loss of self-worth, and little relief; grief, by contrast, typically comes in waves, and positive memories of the deceased usually remain intact
    \item \textbf{Anxiety disorders:} persistent worry, panic, or avoidance that exceeds the normal anxious feelings associated with grief
    \item \textbf{Post-traumatic stress disorder:} following a sudden, violent, or frightening death, with intrusive flashbacks, nightmares, and avoidance of reminders
    \item \textbf{Adjustment disorder:} a less specific emotional reaction to a stressor that does not have the distinctive pattern of grief
    \item \textbf{Medical conditions:} physical causes of fatigue, weight change, or poor sleep, such as thyroid disease, should be considered when symptoms are prominent
\end{itemize}

\section*{When to Seek Medical Care}
Grief itself rarely requires medical treatment, but professional help should be sought when it is severe or prolonged.

\begin{itemize}
    \item Grief is overwhelming and does not begin to ease after many months
    \item The person is unable to carry out daily activities, work, or self-care
    \item Persistent hopelessness, guilt, or thoughts of death or suicide
    \item Heavy drinking, drug misuse, or neglect of physical health as a way of coping
    \item Physical symptoms such as weight loss, persistent chest pain, or poor sleep that cause concern
\end{itemize}

\textbf{Call 000 (or local emergency services) for:}
\begin{itemize}
    \item Thoughts of suicide or plans to harm oneself
    \item Statements that life is not worth living or the person wishes to die
    \item Self-harm or an overdose
    \item Acute confusion or the sudden onset of severe physical symptoms
\end{itemize}

In addition to emergency services, free 24-hour crisis support is available by telephone from services such as Lifeline (13 11 14).

\section*{Diagnosis and Investigations}
\begin{itemize}
    \item \textbf{Clinical interview:} discussion of the loss, the relationship to the deceased, the circumstances and timing of the death, and how the person has been coping
    \item \textbf{Assessment of grief symptoms:} exploring emotions, physical symptoms, daily functioning, and the availability of social support
    \item \textbf{Risk assessment:} careful inquiry into suicidal thoughts, self-harm, and harmful use of alcohol or drugs
    \item \textbf{Screening for co-morbid conditions:} depression, anxiety, and post-traumatic stress disorder frequently coexist with grief and are treatable
    \item \textbf{Physical examination and basic tests:} used to exclude medical causes of fatigue, poor sleep, or weight change where these are prominent
\end{itemize}

\section*{Management}
Most bereaved people need support rather than formal treatment. For those with prolonged or complicated grief, structured help is effective.

\begin{itemize}
    \item \textbf{Social support:} support from family, friends, and community, and the opportunity to talk about the person and the loss, is usually the most important factor in recovery
    \item \textbf{Counselling:} time-limited supportive counselling can help people process the loss and adjust to changed roles
    \item \textbf{Complicated grief therapy:} a specific structured psychological therapy for prolonged grief that helps the person come to terms with the reality of the loss and re-engage with daily life
    \item \textbf{Treatment of co-morbid conditions:} psychological therapy and, where indicated, antidepressant medication for co-existing depression or anxiety
    \item \textbf{Bereavement support groups:} sharing experiences with others who have suffered a similar loss can normalise the process
    \item \textbf{Practical advice:} maintaining routines, postponing major decisions, and limiting alcohol
\end{itemize}

\section*{Complications}
\begin{itemize}
    \item \textbf{Prolonged grief disorder:} severe, persistent grief with major functional impairment beyond 12 months
    \item \textbf{Depression and anxiety disorders:} developing alongside or following the bereavement
    \item \textbf{Post-traumatic stress disorder:} after sudden, violent, or distressing deaths
    \item \textbf{Alcohol and drug misuse:} used as a way of coping with emotional pain
    \item \textbf{Social isolation and relationship difficulties:} including strain on surviving relationships
    \item \textbf{Physical health decline:} from poor sleep, appetite, and immunity
    \item \textbf{Suicidal thoughts and behaviour:} the most serious complication of bereavement
\end{itemize}

\section*{Prognosis and Outcomes}
For most people, acute grief gradually softens over months, and the person learns to live with the loss while still cherishing the relationship with the deceased.

\begin{itemize}
    \item Most bereaved people feel substantially better within one to two years without formal treatment
    \item Grief rarely disappears completely, but it usually becomes less intense and less intrusive over time
    \item Around 1 in 10 people develops prolonged grief disorder, which responds well to specific psychological therapy
    \item Early support and strong social connections are associated with a smoother recovery
    \item Anniversaries and milestones may reactivate grief, which is normal and does not indicate a poor outcome
\end{itemize}

\section*{Prevention and Self-Care}
Grief cannot be prevented, but good self-care reduces the chance of it becoming complicated or prolonged.

\begin{itemize}
    \item Allow yourself to grieve, and do not rush the process or compare it with others
    \item Talk about the person who died and share memories with family and friends
    \item Accept help with practical tasks and company when you need it
    \item Maintain basic routines of sleep, eating, and gentle exercise
    \item Limit alcohol, which can worsen low mood and disrupt sleep
    \item Postpone major decisions, such as moving house or changing jobs, until you feel more settled
    \item Seek professional help early if grief is overwhelming, prolonged, or accompanied by suicidal thoughts
\end{itemize}

\section*{Sources and Further Reading}
The following organisations provide reliable information on bereavement and grief for patients, families, and health professionals.

\begin{itemize}
    \item NHS. \textit{Information on grief, bereavement, and where to find support.} https://www.nhs.uk/
    \item Mayo Clinic. \textit{Overview of grief and coping with the loss of a loved one.} https://www.mayoclinic.org/
    \item World Health Organization. \textit{Resources on mental health and the recognition of prolonged grief disorder.} https://www.who.int/
    \item MSD Manual. \textit{Professional information on grief and bereavement.} https://www.msdmanuals.com/
    \item MedlinePlus. \textit{Health information on grief and bereavement for patients and families.} https://medlineplus.gov/
\end{itemize}
